% Generated draft: requires mathematical and source review.
\documentclass[11pt]{article}
\usepackage[margin=25mm]{geometry}
\usepackage{fontspec}
\setmainfont{Latin Modern Roman}
\usepackage{amsmath,amssymb,mathtools,unicode-math}
\setmathfont{Latin Modern Math}
\usepackage{enumitem,needspace,xurl,hyperref}
\setlength{\parindent}{0pt}
\setlength{\parskip}{5pt}
\setlength{\emergencystretch}{3em}
\newcommand{\N}{\mathbb{N}}
\newcommand{\Z}{\mathbb{Z}}
\newcommand{\Q}{\mathbb{Q}}
\newcommand{\R}{\mathbb{R}}
\newcommand{\C}{\mathbb{C}}
\newcommand{\F}{\mathbb{F}}
\newcommand{\A}{\mathbb{A}}
\newcommand{\PP}{\mathbb{P}}
\newcommand{\E}{\mathbb{E}}
\newcommand{\eps}{\varepsilon}
\newcommand{\dd}{\,\mathrm{d}}
\begin{document}
\section*{UnsolvedMath ID: 30005105}
\textbf{Source problem code:} catalogue entry
\subsection{Definitions and mathematical statement}
The supplied question is: ``Which induced subgraphs must occur in graphs of sufficiently large chromatic number?'' It specifies a research direction rather than a particular conjecture.

For a simple undirected graph \(G=(V,E)\) and \(S\subseteq V\), the induced subgraph \(G[S]\) has vertex set \(S\) and edge set \(\{e\in E:e\subseteq S\}\). The chromatic number \(\chi(G)\) is the least cardinality of a set of colours admitting a vertex colouring in which adjacent vertices receive different colours.

A possible precise framework, conditional on specifying an ambient class \(\mathcal C\) of simple undirected graphs and a family \(\mathcal F\) of candidate induced subgraphs, is the property
\[
\exists k\in\Z_{\geq 1}\;\forall G\in\mathcal C:\quad
\chi(G)\geq k\ \Longrightarrow\
\exists H\in\mathcal F\;\exists S\subseteq V(G)\quad G[S]\cong H.
\]
Here the threshold may depend on \(\mathcal C\) and \(\mathcal F\), but must be independent of \(G\). For a question about one prescribed graph \(H\), one would take \(\mathcal F=\{H\}\).

The source does not specify the graph domain, whether graphs are finite, an ambient class \(\mathcal C\), candidate graphs or families, or whether this particular quantifier pattern is intended. The displayed formula is therefore an explicitly conditional formulation template, not a conjecture attributed to the source. No uniquely determined theorem or conjecture can be recovered from the supplied record.
\subsection{Short English statement}
The question asks what induced structures are forced by sufficiently large chromatic number. The supplied record leaves the graph class, the structures sought, and their quantifiers unspecified, so it supports a research question and a conditional formulation template rather than a definite conjecture.
\subsection{Sources}
\begin{itemize}
\item Supplied local UnsolvedMath registry record, permanent ID \(30005105\), title ``Induced Structures in Highly Chromatic Graphs'', source problem code \texttt{catalogue entry}, in the Oberwolfach Reports Open Problems collection.
\item Supplied report citation: \emph{Combinatorics, Probability and Computing} (2023), Oberwolfach Reports, pp.~1165--1237, DOI \texttt{10.4171/owr/2022/22}. \url{https://doi.org/10.4171/owr/2022/22}. The record gives the proposed year as 2022.
\item The original extracted text names Alex Scott and describes joint work with Ant\'onio Gir\~ao, Freddie Illingworth, Emil Powierski, Michael Savery, Youri Tamitegama and Jane Tan. Its embedded heading reads ``Induced subgraphs of induced subgraphs of large chromatic number''. This attribution does not identify a separate proposer of a precise conjecture.
\item Supplied literature source, labelled ``official OWR report'': \url{https://publications.mfo.de/bitstream/handle/mfo/3964/OWR_2022_22.pdf?isAllowed=y&sequence=4}.
\item Literature assessment dated exactly \texttt{2026-08-22}: the supplied curation labels the problem open because it is a broad methodological question without a uniquely testable claim and with unclear resolution criteria. This is supplied dated curation, not an independently verified literature assessment. The background expressly requests independent verification of time-sensitive status and cited claims.
\item The registry supplies the statement status \texttt{corrected\_verified} and reports that the official report was checked directly and the clean statement preserves the complete mathematical target. These are supplied provenance claims; no independent source inspection or verification is asserted here.
\end{itemize}
\end{document}
